\chapter{三个可交付的毕业项目}
\label{chap:capstone}

\begin{takeaway}[学习目标]
把前面章节组合成可展示的作品。三个项目分别覆盖知识工作、业务自动化和开发者工具；每个项目都有范围、里程碑、数据、评测、安全与演示要求。
\end{takeaway}

\section{作品集不是功能截图}

一个有说服力的 Agent 项目需要回答：用户任务是什么；为什么需要 Agent；系统如何运行；失败如何处理；评测结果怎样；怎样复现。仓库至少包含 README、架构图、最小数据集、评测命令、样例 trace、安全说明和演示视频。

建议统一使用四周节奏：第一周建立任务合同和基线，第二周完成工具与主循环，第三周加入评测和可靠性，第四周进行用户测试与发布。先做到窄而完整，再增加范围。

\section{项目 A：有证据的研究 Agent}

\subsection{任务定义}

输入一个具体研究问题和允许来源，输出结构化报告、事实来源、冲突信息与未知项。适合市场研究、政策整理、技术选型或课程导航。

\subsection{最小范围}

\begin{itemize}
  \item 搜索、打开网页、提取正文、保存证据四个只读工具。
  \item 编排器按子问题委托 3--5 个独立研究任务。
  \item 合并器去重并标记冲突，每个事实必须有 URL 和访问日期。
  \item 生成报告前运行覆盖度与引用支持检查。
\end{itemize}

\subsection{评测集}

准备 30 个问题：10 个单来源事实、10 个多来源比较、5 个没有公开答案、5 个来源相互冲突。指标包括证据召回率、事实正确率、引用精度、拒答质量、P95 时间和单报告成本。

\subsection{安全边界}

不登录私人账户，不下载来源不明的可执行文件，不把网页指令当作任务指令，不自动发布报告。保存最小网页片段并遵守站点访问规则与内容许可。

\subsection{加分项}

提供来源时间线、差异更新、断点恢复、搜索预算和“结论到证据”双向跳转。展示一个失败案例及修复前后的回归数据，比增加华丽 UI 更有价值。

\section{项目 B：客服工单 Agent}

\subsection{任务定义}

读取用户问题与账户的受限信息，检索知识库，生成回复草稿；在授权范围内执行退款申请或升级人工。重点是权限、业务规则与人在回路。

\subsection{最小范围}

\begin{itemize}
  \item 路由：信息问答、账户问题、退款、无法识别。
  \item 工具：查订单、查政策、创建回复草稿、创建退款 proposal、转人工。
  \item 退款金额由代码计算，模型不能直接填写任意金额。
  \item 外发回复和退款提交都使用绑定参数的确认令牌。
\end{itemize}

\subsection{评测集}

从脱敏工单构造 100 条样本，覆盖多语言、情绪表达、缺失订单号、超期退款、重复请求、提示注入和跨账户查询。指标包括路由准确率、政策引用正确率、一次解决率、越权率和人工处理分钟数。

\subsection{可靠性要求}

所有写操作幂等；工具失败后保留草稿；状态变化写审计日志；依赖不可用时降级为人工队列；Agent 不能把“无法查询”编造成账户事实。

\section{项目 C：代码维护 Agent}

\subsection{任务定义}

在隔离工作区内读取仓库、定位问题、修改有限文件、运行测试、审查 diff，最终生成补丁和变更说明，不默认推送远端。

\subsection{最小范围}

\begin{itemize}
  \item 工具：列文件、搜索、读文件、应用补丁、运行允许命令、查看 diff。
  \item 计划中标注将修改的模块和验收命令。
  \item shell 使用命令白名单、工作目录限制、资源与时间上限。
  \item 完成条件必须包含测试通过与 diff 自审。
\end{itemize}

\subsection{评测集}

在固定 commit 上准备 20--50 个真实小任务：修复单元测试、增加校验、更新文档和局部重构。用隐藏测试验证功能，另测改动范围、依赖变化、危险命令和运行成本。可参考 SWE-bench 的任务设计思想，但自己的项目应从小而可复现开始\citep{swebench}。

\subsection{安全边界}

默认无生产凭证和通用网络；不读取工作区外文件；不执行仓库内不可信脚本，除非允许列表明确授权；提交、推送、发布包和部署均需人工确认。

\section{共同交付清单}

\begin{enumerate}
  \item 一条命令启动，一条命令跑评测。
  \item 架构图标出模型、工具、状态、权限和数据边界。
  \item 至少 30 条版本化评测样本和基线结果。
  \item 三条完整 trace：成功、可恢复失败、安全拒绝。
  \item 成本与延迟报告，注明模型和日期。
  \item 威胁模型、已知限制、数据与许可证说明。
  \item 3--5 分钟演示真实任务，不剪掉等待和失败恢复。
\end{enumerate}

\section{项目答辩的六个问题}

\begin{enumerate}
  \item 为什么不用确定性工作流或一次模型调用？
  \item 最常见的三个失败是什么，如何发现？
  \item 哪些动作模型永远无权直接执行？
  \item 当前评测集遗漏了哪些真实分布？
  \item 模型或框架替换后，哪些模块不需要改？
  \item 再给两周，你会根据哪项数据决定下一步？
\end{enumerate}

\begin{practice}[项目立项书]
从三个项目中选一个，写一页立项书：目标用户、任务合同、非目标、架构、风险、四周里程碑、50 条评测样本如何获得、发布门槛。删掉所有无法在四周内验收的功能。
\end{practice}

\begin{checkpoint}
\begin{itemize}
  \item \checkbox 我的项目有真实任务、基线和可复现评测。
  \item \checkbox Demo 同时展示成功、失败恢复和安全拒绝。
  \item \checkbox 仓库公开说明数据来源、许可证和已知限制。
  \item \checkbox 我能用指标而不是感觉解释架构选择。
\end{itemize}
\end{checkpoint}

